\section{Modelagem do banco de dados}\label{sec:modelagem}

A modelagem do banco de dados foi desenvolvida com o objetivo de estruturar as informações utilizadas no processo de avaliação psicopedagógica, mantendo a organização dos registros e a relação entre aprendentes, profissionais, instrumentos avaliativos, aplicações e relatórios. O banco de dados foi implementado em MySQL e organizado segundo o modelo relacional \cite{codd1970_relational}, com utilização de identificadores únicos e chaves estrangeiras para estabelecer as relações entre as entidades. A estrutura foi representada por meio do modelo entidade-relacionamento \cite{chen1976_entity_relationship}, e o acesso aos dados pela aplicação é feito pelo Prisma, ferramenta de mapeamento objeto-relacional, conforme os conceitos apresentados na Seção~\ref{sec:fund_banco}.

O modelo contempla inicialmente a entidade \textit{users}, responsável pelo cadastro dos usuários do sistema. Cada usuário possui identificador único, nome, e-mail, senha armazenada por meio de hash bcrypt com fator de custo 12, perfil de acesso e status. Foram definidos os perfis \textit{ADMIN} e \textit{PSICOPEDAGOGO}, permitindo diferenciar os níveis de acesso às funcionalidades do sistema. A entidade também possui campos para controle de criação, atualização e exclusão lógica dos registros. Associadas aos usuários estão as entidades \textit{refresh\_tokens} e \textit{password\_resets}, utilizadas, respectivamente, para o gerenciamento da renovação das sessões de autenticação e para o processo de recuperação de senha.

A entidade \textit{aprendentes} concentra as informações cadastrais dos indivíduos acompanhados pelo profissional. São armazenados dados como nome, data de nascimento, sexo, informações do responsável, contatos, e-mail e informações relacionadas ao contexto escolar, incluindo escola, série e turno. Também foi disponibilizado um campo para observações. Por envolverem crianças, esses dados estão sujeitos ao art. 14 da LGPD, e as respostas das avaliações e os relatórios, por poderem conter dados de saúde, enquadram-se como dados pessoais sensíveis nos termos do art. 5º, inciso II, da mesma lei. As implicações dessa classificação para o uso em produção são discutidas na Seção~\ref{sec:seguranca}. 

Os campos opcionais da entidade (contato secundário, e-mail, escola, série e turno) permitem ao profissional registrar apenas as informações necessárias ao atendimento, em consonância com o princípio de minimização de dados. A entidade possui o campo \textit{responsavelId}, que estabelece uma chave estrangeira para o usuário responsável pelo acompanhamento daquele aprendente. Esse vínculo é utilizado para implementar a restrição de visibilidade: cada psicopedagogo visualiza somente os aprendentes associados ao seu identificador.

A entidade possui ainda os campos de controle de criação, atualização e exclusão lógica. Para possibilitar o registro longitudinal de informações, a entidade \textit{historico\_aprendentes} mantém registros associados a cada aprendente, contendo uma descrição e uma data, estabelecendo uma relação de um para muitos entre aprendentes e seus históricos.

A representação dos instrumentos utilizados no processo avaliativo foi realizada por meio da entidade \textit{avaliacoes}. Cada avaliação possui nome, descrição, categoria, versão e status, podendo assumir os estados \textit{DRAFT}, \textit{ACTIVE} ou \textit{ARCHIVED}. Um aspecto relevante dessa modelagem é a utilização do campo \textit{parentId}, que estabelece uma autorrelação entre avaliações e permite representar diferentes versões de um mesmo instrumento.

A regra de versionamento é a seguinte: quando uma requisição de edição inclui alteração nas questões de um instrumento que já possui ao menos uma aplicação vinculada, o sistema não modifica o registro existente. Em vez disso, arquiva o registro atual (status \textit{ARCHIVED}) e cria automaticamente um novo registro com versão incrementada, herdando os metadados e recebendo as novas questões. O registro anterior permanece intacto e continua associado a todas as aplicações já realizadas com ele. Quando o instrumento ainda não possui aplicações, as questões são alteradas diretamente no mesmo registro, sem geração de nova versão. 

Alterações apenas em nome, descrição, categoria ou status nunca geram nova versão, independentemente de haver aplicações; essas alterações ficam documentadas no registro de auditoria com os valores anteriores e os novos. Essa regra garante que nenhuma resposta registrada tenha sua questão de origem alterada retroativamente, preservando a rastreabilidade do histórico avaliativo. 

A entidade também possui o campo \textit{autorId}, que referencia o usuário responsável pela criação do instrumento. Esse campo é utilizado para implementar a restrição de visibilidade das avaliações, de modo que cada psicopedagogo acessa somente os instrumentos que criou, seguindo o mesmo princípio adotado para os aprendentes.

Os itens que compõem cada avaliação são armazenados na entidade \textit{questions}. Cada questão está vinculada a uma avaliação específica e possui texto, tipo, obrigatoriedade e posição de apresentação. O modelo contempla diferentes tipos de resposta, incluindo texto curto, texto longo, valor numérico, data, seleção única, seleção múltipla, escala e respostas de sim ou não. Para questões de seleção, são armazenadas as opções correspondentes, enquanto questões de escala podem possuir valores mínimo e máximo e respectivos rótulos. Essa estrutura permite que os instrumentos sejam configurados de maneira flexível, sem que cada protocolo precise possuir uma estrutura de banco de dados específica.

A aplicação de um instrumento a um aprendente é representada pela entidade \textit{aplicacoes}. Essa entidade estabelece a relação entre o aprendente, a avaliação utilizada e o profissional aplicador (\textit{aplicadorId}). Além dessas associações, são registrados o estado da aplicação, a data de aplicação (opcional, pois a aplicação pode ser criada antes do preenchimento efetivo), observações e a data de finalização.

Foram definidos quatro estados, com transições bem delimitadas. A aplicação é criada no estado \textit{PENDING}, indicando que ainda não houve registro de respostas. Ao salvar respostas sem a marcação de rascunho, o estado avança para \textit{IN\_PROGRESS}, sinalizando preenchimento em andamento. Quando o profissional opta explicitamente por salvar como rascunho, o estado passa a \textit{DRAFT}, indicando que o preenchimento está pausado intencionalmente e não deve ser tratado como progresso definitivo. 

A partir de qualquer estado não finalizado, o profissional pode concluir a aplicação, o que a leva ao estado \textit{COMPLETED}; nesse momento, todas as questões obrigatórias são verificadas e as respostas ficam bloqueadas para edição. Aplicações no estado \textit{COMPLETED} não podem ser excluídas. Assim, uma mesma avaliação pode ser aplicada a diferentes aprendentes, enquanto um aprendente pode possuir múltiplas aplicações ao longo de seu acompanhamento.

As respostas fornecidas durante uma aplicação são armazenadas na entidade \textit{respostas}. Cada registro relaciona uma aplicação a uma questão específica e possui um campo \textit{valor} do tipo JSON, permitindo armazenar diferentes formatos de resposta de acordo com o tipo da questão. Foi definida ainda uma restrição de unicidade composta entre \textit{aplicacaoId} e \textit{questionId}, impedindo que uma mesma questão possua mais de uma resposta dentro da mesma aplicação.

Por fim, a entidade \textit{relatorios} representa a devolutiva associada a uma aplicação concluída. Cada relatório está vinculado a uma aplicação específica e possui um autor, além dos campos destinados à conclusão e às observações do profissional. A relação entre \textit{aplicacoes} e \textit{relatorios} é de zero ou um, de modo que uma aplicação pode não possuir relatório inicialmente, mas pode receber uma devolutiva posteriormente. Essa estrutura mantém a elaboração das conclusões vinculada ao registro da avaliação, sem transferir ao sistema a responsabilidade pela interpretação automática dos resultados.

Para a rastreabilidade das operações, o modelo inclui a entidade \textit{audit\_logs}, que registra cada criação, alteração e exclusão de dados e cada consulta a dados clínicos. Cada registro armazena o usuário que realizou a operação (\textit{userId}) e o seu perfil no momento (\textit{userRole}), a ação realizada (\textit{action}: criação, alteração, exclusão, consulta, listagem, login ou tentativa de login recusada), a entidade e o identificador do registro afetado (\textit{entity} e \textit{entityId}), o aprendente a que o dado pertence (\textit{aprendenteId}), o endereço IP de origem, o momento da operação (\textit{timestamp}) e, opcionalmente, as diferenças entre os valores anteriores e os novos (\textit{diff}, do tipo JSON). O campo \textit{aprendenteId} permite reunir todos os acessos aos dados de uma mesma criança, inclusive os realizados por meio de aplicações, respostas e relatórios. A entidade não possui chave estrangeira para \textit{users}, de modo que o registro de auditoria não depende de alterações posteriores no cadastro dos usuários.

De forma geral, os principais relacionamentos do modelo são apresentados no Diagrama Entidade-Relacionamento. Um usuário pode estar associado a várias aplicações e relatórios, um aprendente pode possuir vários registros históricos e aplicações, e uma avaliação pode possuir diversas questões e ser utilizada em diversas aplicações. Cada aplicação, por sua vez, pode possuir várias respostas e, opcionalmente, um relatório. As questões também podem estar relacionadas a várias respostas, considerando diferentes aplicações do mesmo instrumento.

Para favorecer a consulta e o gerenciamento dos dados, foram definidos índices em campos utilizados com maior frequência nas operações do sistema. Entre eles estão o índice único para o e-mail dos usuários e o índice sobre o seu status; os índices sobre o nome, a exclusão lógica e o campo \textit{responsavelId} dos aprendentes; os índices sobre o campo \textit{autorId}, o status, a categoria e a exclusão lógica das avaliações; e os índices sobre os identificadores utilizados nas aplicações e respostas. 

Os índices sobre \textit{responsavelId} e \textit{autorId} são particularmente relevantes para as consultas de listagem, pois essas operações filtram os registros pelo usuário autenticado. A entidade \textit{audit\_logs} possui índices sobre o usuário, sobre a combinação de entidade e identificador do registro, sobre o aprendente e sobre o momento da operação, que correspondem às consultas mais frequentes ao registro de auditoria.

O Diagrama Entidade-Relacionamento desenvolvido para o sistema permite visualizar graficamente as entidades, seus atributos e os relacionamentos estabelecidos entre elas. Em razão da quantidade de entidades, o diagrama foi dividido em duas partes. A Figura~\ref{fig:der_parte1} apresenta as entidades relacionadas aos usuários, à autenticação, aos aprendentes e à auditoria: \textit{users}, \textit{refresh\_tokens}, \textit{password\_resets}, \textit{aprendentes}, \textit{historico\_aprendentes} e \textit{audit\_logs}.

\begin{figure}[H]
    \centering
    \caption{Diagrama Entidade-Relacionamento do sistema - parte 1: usuários, aprendentes e auditoria}
    \label{fig:der_parte1}
    \includegraphics[width=\textwidth,height=0.6\textheight,keepaspectratio]{der_parte1.png}
    \fonte{Elaborado pela autora (2026).}
\end{figure}

A Figura~\ref{fig:der_parte2} apresenta as entidades relacionadas ao processo avaliativo: \textit{avaliacoes}, \textit{questions}, \textit{aplicacoes}, \textit{respostas} e \textit{relatorios}, incluindo a autorrelação de versionamento das avaliações. As entidades \textit{users} e \textit{aprendentes}, detalhadas na primeira parte, aparecem representadas de forma simplificada, em linha tracejada, apenas para indicar seus relacionamentos com as demais entidades.

\begin{figure}[H]
    \centering
    \caption{Diagrama Entidade-Relacionamento do sistema - parte 2: instrumentos, aplicações e relatórios}
    \label{fig:der_parte2}
    \includegraphics[width=\textwidth,height=0.6\textheight,keepaspectratio]{der_parte2.png}
    \fonte{Elaborado pela autora (2026).}
\end{figure}
